\documentclass[12pt]{article}

\usepackage{amsmath, amssymb, amsthm}
\usepackage{enumerate}
\usepackage{hyperref}
\usepackage{geometry}
\usepackage{newtxtext}
\usepackage{newtxmath}

\geometry{margin=1in}

\title{Probability Theory: Lecture 1}
\author{}
\date{}

% ----------------------------------------------------------------
\begin{document}
\maketitle

We present a classic problem with a very conceptual 
solution, showcasing linearity of expectation. 
The problem is known as \emph{Buffon's 
needle}, and one elegant solution strategy is due to 
Barbier.

We follow Barbier's exposition step by step, beginning with Buffon's needle and arriving, through successive generalisations, at Buffon's noodle. The key
insight is to use probabilistic reasoning, specifically linearity of
expectation.

Please do not worry about rigour and focus instead on the key ideas. Please do not feel guilty about skipping the optional parts, which may be essential to the solution but not necessary to appreciate the ideas of finding the right abstraction.

% ================================================================
\section*{Setup}

Parallel straight lines are drawn on the plane $\mathbb{R}^2$ at equal
distance $d$ apart. A needle of length $L$ that is less than or equal to $d$ is thrown at random onto the
plane. We assume that all positions and orientations of the needle are
equally likely. \textbf{What is the probability that the needle crosses at
least one of the lines?}

% ================================================================
\section*{Questions}

\begin{enumerate}[(a)]

  \item \textbf{Long needles and counting crossings.}
  Instead of asking for a probability directly, we work with the random
  variable $X_1 =$ the number of times a randomly thrown needle of length
  $L_1$ crosses one of the parallel lines. (A needle can cross more than one
  line if it is long enough.)

  \begin{enumerate}[(i)]
    \item We say the needle is \emph{short} if $L_1 < d$. What integer
      values can $X_1$ take when the needle is short? Briefly justify
      your answer.

    \item What integer values can $X_1$ take when the needle can be
      arbitrarily long?

    \item Let $p_n = P(X_1 = n)$ for $n = 0, 1, 2, \ldots$\,. Write
      down the definition of $E[X_1]$ as a sum involving the $p_n$.

    \item Show that when the needle is short, $E[X_1] = p_1$, the
      probability of crossing exactly one line. Why does this mean it
      suffices to compute $E[X_1]$ in order to solve Buffon's original
      problem?

    \item Explain why $E[X_1]$ depends only on $L_1$ (and on the fixed
      spacing $d$), so that we may write $E[X_1] = f(L_1)$ for some
      function $f$.
  \end{enumerate}

  \item \textbf{Two needles and linearity of expectation.}
  Throw a second needle of length $L_2$ independently, and let $X_2$ be
  the number of times it crosses the lines.  We now weld the two needles
  rigidly together at one point to form a single rigid object, and throw
  the combined object at random.

  Recall that for \emph{any} two random variables $X$ and $Y$ (with finite expected values):
  \[
    E[X + Y] = E[X] + E[Y].
  \]
  Independence is \textbf{not} required for this identity.

  \begin{enumerate}[(i)]
    \item (Optional) Argue that $X_1$ and $X_2$ are independent
      when the two needles are thrown separately.

    \item Using linearity of expectation, write $E[X_1 + X_2]$ in terms
      of $E[X_1]$ and $E[X_2]$.

    \item Use part~(a)(v) to write $E[X_1 + X_2]$ as a function of
      $L_1$ and $L_2$.
  \end{enumerate}

  \item \textbf{Many needles and the functional equation.}
  We now weld $n$ needles of lengths $L_1, \ldots, L_n$ together rigidly
  (end-to-end, or in any rigid configuration) to form a polygonal
  \emph{noodle}, and throw the whole thing at random.

  \begin{enumerate}[(i)]
    \item Write down $E[X_1 + X_2 + \cdots + X_n]$ using linearity of
      expectation.

    \item Use part~(a)(v) to express $E[X_1 + \cdots + X_n]$ purely in
      terms of the function $f$ and the lengths $L_1, \ldots, L_n$.

    \item (Optional) The welded object of total length $L_1 + \cdots + L_n$ is
      itself a single rigid object, so by~(a)(v) its expected crossing
      number is $f(L_1 + \cdots + L_n)$. Deduce that
      \[
        f(L_1 + \cdots + L_n) = f(L_1) + \cdots + f(L_n).
      \]
      Use this to prove that $f(qL) = q\,f(L)$ for every positive
      rational $q$, and hence that $f$ is linear on the positive
      rationals: $f(L) = rL$ for some constant $r > 0$ and all rational
      $L > 0$.

    \item (Optional) Explain why $f$ is monotonically increasing, and use this to
      conclude that $f(L) = rL$ for \emph{all} real $L > 0$.
  \end{enumerate}

  \item \textbf{(Optional) Buffon's noodle: a smooth curve.}
  Now let $C$ be any smooth rigid curve of length $L$, and let $Y$ be the
  number of times $C$ crosses the parallel lines when thrown at random.

  \begin{enumerate}[(i)]
    \item (Optional) Approximate $C$ by a polygonal path with $k$ straight segments
      of lengths $\ell_1, \ldots, \ell_k$ (so $\ell_1 + \cdots + \ell_k
      \approx L$). Using part~(c) and a limiting argument, explain why
      $E[Y] = rL$.

    \item (Optional) Conclude that $E[Y] = rL$ holds for any rectifiable curve $C$
      of length $L$, for the same constant $r$ as before (which we have
      not yet determined).
  \end{enumerate}

\newpage

  \item \textbf{Determining $r$: the circle trick.}
  
  The main result to use is the formula $E[Y] = rL$ for any curve $C$ of length $L$.
  It is OK if you have skipped the earlier exercises.

  We now choose a very special curve $C$ to pin down the constant $r$.
  Let $C$ be a \emph{circle of diameter $d$} (so it has the same diameter
  as the spacing between the lines).

  \begin{enumerate}[(i)]
    \item By thinking carefully about the geometry, compute $E[Y]$
      directly from the definition of expectation. (Hint: how many times
      does a circle of diameter $d$ cross a family of parallel lines
      spaced $d$ apart, regardless of where it lands?) You should
      obtain $E[Y] = 2$.

    \item Compute the circumference $L$ of this circle.

    \item Use $E[Y] = rL$ to solve for $r$.

    \item Finally, for a \emph{short} needle of length $L_1 < d$, use
      part~(a)(iv) to deduce
      \[
        P(\text{needle crosses a line}) \;=\; \frac{2L_1}{\pi d}.
      \]
  \end{enumerate}

\end{enumerate}

% ================================================================
\newpage
\section*{Solutions}

\begin{enumerate}[(a)]

  % ------------------------------------------------------------------
  \item

  \begin{enumerate}[(i)]

    \item A short needle has length $L_1 < d$, the spacing between lines.
    Since consecutive lines are distance $d$ apart and the needle is
    strictly shorter than $d$, it cannot reach from one line to the next,
    so it can cross \emph{at most} one line. Therefore $X_1 \in \{0, 1\}$.

    \item When $L_1$ can be arbitrarily large there is no upper bound on
    the number of lines it can cross, so $X_1$ can take any value in
    $\{0, 1, 2, 3, \ldots\}$.

    \item By the definition of expectation for a non-negative integer-valued
    random variable,
    \[
      E[X_1] \;=\; \sum_{n=0}^{\infty} n\, p_n
              \;=\; 0 \cdot p_0 + 1 \cdot p_1 + 2 \cdot p_2 + \cdots
    \]

    \item When $L_1 < d$ we have $p_n = 0$ for all $n \geq 2$ (by
    part~(i)), so
    \[
      E[X_1] = 0 \cdot p_0 + 1 \cdot p_1 = p_1.
    \]
    But $p_1 = P(X_1 = 1) = P(\text{needle crosses exactly one line}) =
    P(\text{needle crosses at least one line})$, which is exactly the
    probability we want to find. So computing $E[X_1]$ is equivalent to
    solving Buffon's original problem.

    \item The parallel lines are uniformly spaced and the needle is thrown
    uniformly at random (all positions and orientations equally likely).
    There is no preferred position or direction on the plane, so the
    distribution of $X_1$ is completely determined by the geometry: only
    the length $L_1$ of the needle and the spacing $d$ matter.  Since $d$
    is fixed throughout, $E[X_1]$ depends only on $L_1$, and we may write
    $E[X_1] = f(L_1)$.

  \end{enumerate}

  % ------------------------------------------------------------------
  \item

  \begin{enumerate}[(i)]

    \item When the two needles are thrown
    separately, the position and orientation of each needle are chosen
    independently of the other. Therefore the events involving the first
    needle are independent of those involving the second, so $X_1$ and
    $X_2$ are independent.

    \item After welding, the total number of crossings of the combined
    object is $X_1 + X_2$ (each needle still contributes its own
    crossings). By linearity of expectation,
    \[
      E[X_1 + X_2] = E[X_1] + E[X_2].
    \]

    \item By part~(a)(v), $E[X_1] = f(L_1)$ and $E[X_2] = f(L_2)$, so
    \[
      E[X_1 + X_2] = f(L_1) + f(L_2).
    \]

  \end{enumerate}

  % ------------------------------------------------------------------
  \item

  \begin{enumerate}[(i)]

    \item Applying linearity of expectation to $n$ needles,
    \[
      E[X_1 + X_2 + \cdots + X_n] = E[X_1] + E[X_2] + \cdots + E[X_n].
    \]

    \item Using $E[X_i] = f(L_i)$ for each $i$,
    \[
      E[X_1 + \cdots + X_n] = f(L_1) + f(L_2) + \cdots + f(L_n).
    \]

    \item The welded object is a single rigid piece of total length
    $L_1 + \cdots + L_n$, so by part~(a)(v) its expected crossing number
    is also $f(L_1 + \cdots + L_n)$.  Combining with part~(ii) above:
    \[
      f(L_1 + \cdots + L_n) = f(L_1) + \cdots + f(L_n). \tag{$*$}
    \]
    Setting all $L_i = L$ in $(*)$ gives $f(nL) = n\,f(L)$ for every
    positive integer $n$.  Setting $L \mapsto L/m$ then gives
    $f(L) = m\,f(L/m)$, i.e.\ $f(L/m) = f(L)/m$.  Combining:
    $f((n/m)L) = (n/m)\,f(L)$ for all positive integers $n, m$.
    Taking $L = 1$ shows $f(q) = q\,f(1)$ for every positive rational
    $q$, i.e.\ $f(L) = rL$ with $r = f(1) > 0$.

    \item If $L > L'$ then a needle of length $L$ can always land in
    every position a needle of length $L'$ can, and sometimes cross more
    lines, so $f(L) \geq f(L') $; in fact $f$ is strictly increasing.
    For any real $L > 0$, choose rationals $q^- < L < q^+$; then
    $rq^- = f(q^-) \leq f(L) \leq f(q^+) = rq^+$.  Sending $q^\pm \to L$
    gives $f(L) = rL$ for all real $L > 0$.

  \end{enumerate}

  % ------------------------------------------------------------------
  \item

  \begin{enumerate}[(i)]

    \item Approximate $C$ by a polygonal path with $k$ segments of lengths
    $\ell_1, \ldots, \ell_k$.  Each segment is a straight needle; call
    $X^{(j)}$ the number of crossings of segment~$j$.  By part~(c),
    \[
      E\!\left[\sum_{j=1}^k X^{(j)}\right]
        = \sum_{j=1}^k f(\ell_j) = r\sum_{j=1}^k \ell_j.
    \]
    As the mesh of the polygonal approximation tends to zero, $\sum \ell_j
    \to L$ and the expected crossings of the polygon converge to $E[Y]$,
    giving $E[Y] = rL$.

    \item The argument holds for any rectifiable curve $C$: approximate by
    polygons and pass to the limit.  The constant $r$ is the same because
    it was determined purely by the spacing $d$ and does not depend on the
    shape of $C$.

  \end{enumerate}

  % ------------------------------------------------------------------
  \item

  \begin{enumerate}[(i)]

    \item A circle of diameter $d$ has the same width as the spacing
    between lines.  No matter where or how the circle lands, it always
    crosses \emph{exactly} two of the parallel lines (the leftmost and
    rightmost lines it straddles), never more and never fewer.  (It lies
    strictly between two consecutive lines only if its centre is more than
    $d/2$ from each line, which is impossible when the diameter equals $d$.)
    Therefore $Y = 2$ with probability $1$, and
    \[
      E[Y] = 2.
    \]
    This is why a circle of diameter $d$ is such a convenient choice: its
    crossing number is \emph{deterministic}, making $E[Y]$ trivial to
    compute.

    \item The circumference of a circle of diameter $d$ is
    \[
      L = \pi d.
    \]

    \item From $E[Y] = rL$ we get $2 = r \cdot \pi d$, so
    \[
      r = \frac{2}{\pi d}.
    \]

    \item For a short needle of length $L_1 < d$, parts~(a)(iv) and the
    formula $f(L_1) = rL_1$ give
    \[
      P(\text{needle crosses a line})
        = E[X_1]
        = f(L_1)
        = \frac{2L_1}{\pi d}.
    \]
    This is the answer to Buffon's original problem, obtained without
    evaluating any trigonometric integral.

  \end{enumerate}

\end{enumerate}

\end{document}